\begin{abstract}
We refine the cyclic construction in OpenAI Math family 160,
\emph{Quantitative Superexponential Bounds for van der Waerden Numbers}.
For absolute constants \(C,K>0\) and every integer \(k\ge K\), we prove
\[
 \log W_2(k)>
 \frac{k}{24}\bigl(\log k-5\log\log k\bigr)-Ck .
\]
The construction supplies more than nonmonochromaticity: every \(k\)-term
progression differs from every constant coloring in at least
\(\fl{b k^{2/3}(\log k)^{5/3}}\) positions, for an absolute \(b>0\).
A direct three-color version has leading coefficient \(1/18\).
Combining the two cyclic constructions gives a uniform all-color
bound with an exactly optimized binary--ternary coefficient.
The improvements come from a two-point norm-band bound, balancing
the light positions left outside heavy blocks, an elementary estimate
for the small-prime dilation, and a joint optimization of two distinct
signature costs.
We give the entire ordinary mathematical argument, including the
geometric and probabilistic ingredients inherited from the source.
This is a quantitative strengthening of that repository result, not a
new claim to have first resolved the underlying superexponential-growth
problem. No Lean formalization of the strengthened statements is claimed.
\end{abstract}

\section{Results, provenance, and the selected research problem}
\label{sec:results}

\subsection{What is being estimated}

Write \(W_r(k)\) for the least integer \(n\) such that every coloring of
\([n]=\{1,\ldots,n\}\) with at most \(r\) colors has a monochromatic
arithmetic progression of length \(k\) and positive common difference.
Unless a base is displayed, logarithms are natural. All asymptotic
statements below concern \(k\to\infty\).

For a coloring \(c\) and a \(k\)-term progression \(P=(x_j)_{j<k}\),
define its distance from monochromaticity by
\begin{equation}\label{eq:distance}
 d_{\rm mono}(c,P)
 =\min_b \#\{j<k:c(x_j)\ne b\}
 =k-\max_b\#\{j<k:c(x_j)=b\}.
\end{equation}
Thus \(d_{\rm mono}(c,P)\ge t\) says that at least \(t\) positions
must be recolored to make \(P\) monochromatic. For two colors, it says
that each color appears at least \(t\) times. For more colors it is a
statement about distance from \emph{each constant word}; it does not
say that every available color occurs.

\begin{theorem}[Two-color bound with robustness]\label{thm:main-two}
There are absolute constants \(C,b>0\) and an integer \(K\) such that,
for every integer \(k\ge K\), there is a cyclic group \(\Z/N_2\Z\) and
a two-coloring \(c\) satisfying
\begin{align}
 \log N_2&\ge
 \frac{k}{24}\bigl(\log k-5\log\log k\bigr)-Ck,
 \label{eq:main-two-size}\\
 d_{\rm mono}(c,P)&\ge
 \fl{b k^{2/3}(\log k)^{5/3}}
 \label{eq:main-distance}
\end{align}
for every nonzero-step \(k\)-term cyclic progression \(P\).
The same robustness holds on an integer interval of length
\((k-1)N_2\). In particular,
\[
 W_2(k)>(k-1)
 \exp\!\left\{\frac{k}{24}
       \bigl(\log k-5\log\log k\bigr)-Ck\right\},
 \qquad
 \liminf_{k\to\infty}\frac{\log W_2(k)}{k\log k}\ge\frac1{24}.
\]
\end{theorem}

The loss \(5k\log\log k/24\) is part of the proved bound. Its optimality
is not asserted. Neither \(C\) nor the numerical onset \(K\) is optimized;
the proof fixes them by explicit inequalities between absolute
constants and elementary functions.

\begin{theorem}[Uniform palette improvement]\label{thm:main-many}
For an integer \(r\ge2\), set
\begin{equation}\label{eq:palette-coefficient}
 m=\fl{\log_2 r},\qquad
 e(r)=
 \begin{cases}
 1,&r\ge3\cdot2^{m-1},\\
 0,&r<3\cdot2^{m-1},
 \end{cases}
 \qquad
 \Lambda(r)=\frac{m}{24}+\frac{e(r)}{72}.
\end{equation}
There are absolute \(C,b,K>0\), independent of \(r\), such that
for all integers \(k\ge K\) and \(r\ge2\),
\begin{equation}\label{eq:main-many}
 W_r(k)>(k-1)\exp\!\left\{
 \Lambda(r)k\bigl(\log k-5\log\log k\bigr)-Cmk
 \right\}.
\end{equation}
There is an \(r\)-coloring of an interval at least this long for which
every \(k\)-term progression has distance at least the quantity in
\eqref{eq:main-distance} from every monochromatic word.
\end{theorem}

The coefficient has a small but exact palette effect:
\begin{center}
\begin{tabular}{@{}c c c@{}}
\toprule
Available colors \(r\)&Digits used&\(\Lambda(r)\)\\
\midrule
2&one binary&\(1/24\)\\
3&one ternary&\(1/18\)\\
4--5&two binary&\(1/12\)\\
6--7&one binary, one ternary&\(7/72\)\\
8--11&three binary&\(1/8\)\\
12--15&two binary, one ternary&\(5/36\)\\
\bottomrule
\end{tabular}
\end{center}
The uniformity in \(r\) comes from constructing just two fixed
palettes before taking digit products.

\subsection{Repository review and choice of direction}

The review was anchored to the public OpenAI release commit
\href{https://github.com/openai/math/commit/adc7f1241b42e322a6451854ab7e4b4c146bf78a}{\texttt{adc7f1241b42}}
(6 October 2026) and the ProveIt checkpoint
\href{https://github.com/VladimirReshetnikov/ProveIt/commit/bbde2b873b2dc7b48af1730926452bdda18bf9c8}{\texttt{bbde2b873b2d}}
(7 October 2026). The source inventory includes the full commit identifiers.
The OpenAI catalog describes 722 manuscripts in 372 families and
explicitly distinguishes different verification stages
\cite{openaiReadme,openaiContents}.
The review here is a targeted comparison of the arithmetic-progression
families and the active ProveIt Ramsey work, rather than a certification
of every result in either repository.

\begin{center}
\small
\begin{tabular}{@{}>{\raggedright\arraybackslash}p{.19\textwidth}>{\raggedright\arraybackslash}p{.35\textwidth}>{\raggedright\arraybackslash}p{.37\textwidth}@{}}
\toprule
Source&Relevant result or status&Consequence for this project\\
\midrule
OpenAI family 159&
Claims \(r_k(N)\le C_kN\) \(\exp[-c_k(\log N)^{\varepsilon_k}]\);
the described comparator emphasizes the reciprocal-sum consequence
\cite{openai159,openai159Scope}.&
A strong density-theorem overlap, but its paper already contains
several natural weighted and coloring consequences.\\
\addlinespace
OpenAI family 160&
Proves the displayed bound
\(W_r(k)>k^{10^{-5}k\lfloor\log_2r\rfloor}\);
the scope file records this formalized endpoint
\cite{openai160,openai160Scope}.&
An elementary, sufficiently explicit construction whose quantitative
losses can be inspected and improved.\\
\addlinespace
ProveIt Ramsey&
The inspected checkpoint records 113 of 120 exact Gowers statements,
with seven open; recent intake includes sharp hereditary-energy
results and failures of universal lifting interfaces
\cite{proveitStatus,proveitIntake}.&
Motivates exact bookkeeping, careful quantifiers, and a complete
proof boundary. A lower-bound project avoids duplicating the large
existing collection of local upper-bound refinements.\\
\bottomrule
\end{tabular}
\end{center}

The same ProveIt status file records 137 vendored OpenAI Lean modules
in its build, an exact qualitative Szemer\'edi statement proved by a
route distinct from the original Gowers argument, and completed
Freiman and Balog--Szemer\'edi interfaces using OpenAI modules
\cite{proveitStatus}. Thus the intersection is already concrete:
arithmetic-progression research is being assembled with explicit
theorem interfaces and a record of which claims are checked.
These are statements of the inspected repository checkpoint; this
report did not independently run that Lean build.

The direct research question selected was:
\emph{how far can the coefficient in family 160 be increased while
retaining a fully justified growing-dimensional construction, and can
the resulting colorings be made quantitatively resistant to
monochromaticity?}
The answer is \cref{thm:main-two,thm:main-many}.

\subsection{The precise increment over prior work}

The underlying cyclic embedding, adaptive mesh, anchored counting
argument, local lemma, and norm-indexed perturbation are due to
OpenAI's source manuscript \cite{openai160}.
They are reconstructed below with the changed hypotheses visible.
The new quantitative steps established in this report are:
\begin{enumerate}
 \item A norm-band key has at most \emph{two} occurrences on a
 progression, improving the source's upper bound of four.
 \item The local lemma simultaneously balances residual light sets
 after heavy blocks are removed. This yields a
 \((1/2-o(1))\)-balanced-or-affine dichotomy in two colors; balancing
 the entire light set alone would not justify this conclusion.
 \item The short-period cutoff can be chosen of order
 \(D(\log k)^3\) at balance accuracy \(1/\log k\), and
 \(\log\lambda=O(D(\log k)^3)\) follows from an elementary prime-count
 estimate. This replaces the much larger power-scale cutoff in
 the source and avoids an unnecessary logarithmic loss.
 \item Global label entropy and quadratic-signature entropy are
 optimized together. Their balance gives
 \(D\asymp k^{1/3}(\log k)^{-2/3}\) and
 \(p\asymp k^{-1/3}(\log k)^{5/3}\), and hence the coefficient \(1/24\).
 \item A bounded-variable lower-tail argument and an exception-set
 count give the robustness estimate. A three-color construction and
 an optimized mixed-radix product give \(\Lambda(r)\).
\end{enumerate}
The finite periodic extension factor \(k-1\) is an existing observation,
used explicitly by Campos--Fox--Schildkraut
\cite[Lemma 2]{camposFoxSchildkraut2026}; it is credited and proved
again in the robust setting.

These are quantitative improvements to a result already present in
the release, rather than a second claim of priority for
superexponential growth. A bounded search of the inspected ProveIt
Ramsey material and then-current incoming reports found no report
containing these refinements. That observation is not a proof of
global bibliographic novelty.

\subsection{Comparison with the contemporary literature}

Fox and Hunter prove three-color superexponential growth, including
\(W_3(k)>2^{k\log^*k/4}\) for large \(k\)
\cite[Theorem 1]{foxHunter2026}.
Campos, Fox, and Schildkraut prove
\(W_2(k)\ge(1-o(1))k2^{k-1}\)
\cite[Theorem 1]{camposFoxSchildkraut2026}.
The source paper \cite{openai160} subsequently supplies the
two-color \(k^{ck}\) construction used here.
The displayed two-color leading coefficient in this report is
\((1/24)/10^{-5}=12500/3\) times the displayed source coefficient.
This compares asymptotic exponents, not practical lower bounds for
small numerical values of \(k\).

Our uniform all-color bound is useful down to two and three colors.
It is not the strongest known estimate in every growing-color regime.
For each fixed \(0<\varepsilon<1\), Fox--Hunter
\cite[Theorem 2]{foxHunter2026} prove
\[
 W_r(k)\ge r^{(1-\varepsilon)k\log k}
 \quad\text{when }\
 k\ge k_0(\varepsilon),\qquad
 r\ge(\log k)^{3/\varepsilon}.
\]
For example, their choice \(\varepsilon=1/2\) is stronger than our
leading coefficient once its hypotheses apply.
Both bounds can simply be retained in the appropriate regime.

\subsection{Proof conventions and status}

All probabilistic spaces used below are finite. Equality at every mesh
boundary is handled by a fixed half-open convention; key identities
retain their actual labels and are never renamed when events are
identified. Constants denoted \(C_0,C_1,\ldots\) are absolute, and
subscripts such as \(C_A\) indicate dependence on a previously fixed
absolute choice. The choice order is
\[
  \text{geometric constants}\quad\longrightarrow\quad A
  \quad\longrightarrow\quad B\quad\longrightarrow\quad C
  \quad\longrightarrow\quad K .
\]
The proof of the strengthened results is an ordinary mathematical
proof presented in this report. Source scope files are not a substitute
for running a prover; no build of the OpenAI or ProveIt Lean library
was performed here. The exact-arithmetic companion checks finite
identities and the palette optimization, and is not represented as a
certificate of the asymptotic existence theorem.
